\section{The divergence bridge in full}

This chapter gives the complete divergence bridge as one continuous
derivation: from a positive cycle to a contradiction, and then from
the absence of cycles to unboundedness.

\subsection{Assumption}

Assume
\begin{equation*}
  \OrbitCycle(7).
\end{equation*}
The goal of the bridge is a contradiction.  Once the contradiction is
established, the general fact
\begin{equation*}
  \neg\,\OrbitCycle(7)\ \Longrightarrow\ \mathrm{IsUnboundedOrbit}(7)
\end{equation*}
finishes the proof.

\subsection{Extraction of the cycle word}

There are indices $m<n$ with
\begin{equation*}
  \Orbit_7(m)=\Orbit_7(n).
\end{equation*}
Put $c=m$, $p=n-m$, and
\begin{equation*}
  w=\mathrm{cycleWord}(c,p).
\end{equation*}
The cycle word has length $p$ and is valid from $m_0=\Orbit_7(c)$,
and every entry is the exact step weight at the corresponding depth.
The extraction is compiled in
\texttt{orbit\_\allowbreak cycle\_\allowbreak imp\_\allowbreak
cycle\_\allowbreak qb8\_\allowbreak input}.

\subsection{The closedness equation}

Closedness gives
\begin{equation*}
  \wordOrbit(w,m_0)=m_0.
\end{equation*}
The word-orbit identity gives the cycle equation
\begin{equation*}
  2^T\,m_0=5^p\,m_0+\wordA(w),
  \qquad T=\wordWeight(w).
\end{equation*}
In particular $T>p\log_2 5$.

\subsection{The QB-8 split}

Filter the word into rise entries and C3 entries:
\begin{equation*}
  \mathrm{rise}=\{t\in w:t<3\},
  \qquad
  \mathrm{c3}=\{t\in w:t\ge3\}.
\end{equation*}
Both filters are nonempty:
\begin{itemize}[leftmargin=*]
\item if $\mathrm{c3}$ were empty, then $T\le2p$, contradicting
  $T>p\log_2 5>\allowbreak2p$;
\item if $\mathrm{rise}$ were empty, then every step would satisfy
  $(5x+1)/2^t<x$, so the states could not close.
\end{itemize}
The split is compiled in \texttt{CycleQb8Input}.

\subsection{The reset block}

Choose a rise start.  The compiled extraction produces parameters
$j,k_0,t,\delta,s$ and states $r,r_j$ such that
\begin{itemize}[leftmargin=*]
\item $t\in\{1,2\}$, $\delta=1$ for $t=1$,
  $\delta\in\{1,3\}$ for $t=2$;
\item $s5^{k_0}=r+1$;
\item $\GeneralOrbit(r)$ and $\FullOrbit(r_j)$;
\item $\mathrm{ResetHeadEq}(s,j,k_0,t,\delta,r_j)$;
\item $s<5^{j-1-k_0}$, $s$ odd, $5\nmid s$.
\end{itemize}
This is the corrected reachability predicate
$\ResetWindow(j,k_0,t,\delta,s)$.

\subsection{The block tail and the unified core}

The block head $r_j$ is followed by a tail of length $L$ ending at
$r_s$.  With $n=s-j$ and $\Delta=W_s-W_j$,
\begin{equation*}
  2^{\Delta}\,r_s=5^n\,r_{j,0}+B.
\end{equation*}
The unified core applies to this tail.  When its final inequality is
closed, the projection interface
\begin{verbatim}
decisiveWindowValuationBoundCorrected_of_unified_core
\end{verbatim}
produces the corrected window bound.  That interface is compiled but
consumes \texttt{unifiedCoreClosed}, whose Lean representative
\texttt{unified\_\allowbreak core\_\allowbreak final\_\allowbreak
no\_\allowbreak hge} is pending.

\subsection{The valuation bounds}

For $t=2$, the block-layer balance gives
\begin{equation*}
  \vtwo{5^{k_0+1}s+\delta5^j}\ge2j+11,
\end{equation*}
and the corrected window bound gives
\begin{equation*}
  \vtwo{5^{k_0+1}s+\delta5^j}\le2j+10.
\end{equation*}
For $t=1$,
\begin{equation*}
  \vtwo{5^{k_0+1}s+5^j-2}\ge2j+12,
\end{equation*}
and
\begin{equation*}
  \vtwo{5^{k_0+1}s+5^j-2}\le2j+11.
\end{equation*}
Both cases are contradictions.

\subsection{The failure window}

The construction assembles
\begin{equation*}
  \mathrm{FailureWindow}(m_0,S,p,w,\mathrm{rise},\mathrm{c3},
  j,k_0,t,\delta,s),
\end{equation*}
whose fields are:
\begin{itemize}[leftmargin=*]
\item \texttt{hj\_lt}: $j<p$;
\item \texttt{ht}, \texttt{hdelta}: the reset weight and parameter;
\item \texttt{hreset}: the exact reset equation;
\item \texttt{hreach}: the corrected reachability predicate;
\item \texttt{hs\_odd}, \texttt{hs\_not\_five}, \texttt{hk},
  \texttt{hs\_lt}: parity, coprimality, and size;
\item \texttt{hfail\_t1}, \texttt{hfail\_t2}: the lower bounds.
\end{itemize}
The existence statement
\texttt{failureWindowExistence} is pending.

\subsection{The contradiction}

Given a failure window and the corrected bound, the packaged lemma
\texttt{failure\_\allowbreak window\_\allowbreak contradicts\_\allowbreak
window\_\allowbreak corrected} derives \texttt{False}.  This is
compiled.

\subsection{No cycle implies unbounded}

The general lemma
\begin{verbatim}
unbounded_of_no_cycle (x : Nat) (h : not OrbitCycle x) :
  IsUnboundedOrbit x
\end{verbatim}
is compiled.  Its proof uses the pigeonhole principle on the finite
set of bounded values.

\subsection{The full chain}

\begin{verbatim}
OrbitCycle 7
  -> cycleWord extraction (compiled)
  -> QB-8 split (compiled)
  -> reset block and reachability (compiled)
  -> unified core (unified_core_final_no_hge pending)
  -> corrected window bound (compiled conditional)
  -> lower bound from rise_block_balance (compiled)
  -> FailureWindow existence (pending)
  -> failure_window_contradicts_window_corrected (compiled)
  -> not OrbitCycle 7
  -> IsUnboundedOrbit 7 (assembly target)
\end{verbatim}

\subsection{Status}

\begin{center}
\small
\begin{tabular}{@{}p{0.62\textwidth}p{0.32\textwidth}@{}}
\toprule
Step & Status \\
\midrule
cycle-word extraction & compiled \\
QB-8 split & compiled \\
reset block and reachability & compiled \\
block-layer lower bounds & compiled \\
corrected window projection & compiled (conditional) \\
failure-window existence & pending \\
failure-window contradiction & compiled \\
no-cycle-to-unbounded lemma & compiled \\
final theorem assembly & pending \\
\bottomrule
\end{tabular}
\end{center}
